% Generated draft: requires mathematical and source review.
\documentclass[11pt]{article}
\usepackage[margin=25mm]{geometry}
\usepackage{fontspec}
\setmainfont{Latin Modern Roman}
\usepackage{amsmath,amssymb,mathtools,unicode-math}
\setmathfont{Latin Modern Math}
\usepackage{enumitem,needspace,xurl,hyperref}
\setlength{\parindent}{0pt}
\setlength{\parskip}{5pt}
\setlength{\emergencystretch}{3em}
\newcommand{\N}{\mathbb{N}}
\newcommand{\Z}{\mathbb{Z}}
\newcommand{\Q}{\mathbb{Q}}
\newcommand{\R}{\mathbb{R}}
\newcommand{\C}{\mathbb{C}}
\newcommand{\F}{\mathbb{F}}
\newcommand{\A}{\mathbb{A}}
\newcommand{\PP}{\mathbb{P}}
\newcommand{\E}{\mathbb{E}}
\newcommand{\eps}{\varepsilon}
\newcommand{\dd}{\,\mathrm{d}}
\begin{document}
\section*{UnsolvedMath ID: 30001310}
\textbf{Source problem code:} catalogue entry
\subsection{Definitions and mathematical statement}
The supplied question is: ``Can the canonical projective structure on a hyperbolic three-manifold be nontrivially deformed?''

A hyperbolic three-manifold is a smooth three-dimensional manifold \(M\) with a Riemannian metric of constant sectional curvature \(-1\). The record does not specify completeness, finite volume, compactness, orientability, or boundary conditions.

Under the conventional interpretation of ``projective'' as \emph{real projective}, a projective structure on \(M\) is a maximal atlas of charts into \(\R\mathrm{P}^{3}\) whose transition maps are restrictions of elements of \(\mathrm{PGL}(4,\R)\). The Klein model identifies hyperbolic three-space with the interior of a projective ball. Its isometries act by projective transformations, so the hyperbolic charts induce a canonical real projective structure, denoted \(\mathcal P_{\mathrm{hyp}}\).

One precise interpretation of local nontrivial deformation uses marked projective structures: two structures on the same smooth manifold are equivalent if a diffeomorphism isotopic to the identity pulls one back to the other. With this convention, the deformation question for a specified \(M\) asks whether there exist \(\varepsilon>0\) and a smooth family \((\mathcal P_t)_{t\in(-\varepsilon,\varepsilon)}\) of real projective structures on \(M\) such that
\[
\mathcal P_0=\mathcal P_{\mathrm{hyp}},
\qquad
\forall\delta\in(0,\varepsilon)\ \exists t\in(-\delta,\delta)
\quad [\mathcal P_t]\ne[\mathcal P_{\mathrm{hyp}}].
\]
Here brackets denote equivalence under the stated marking convention, and smoothness means smooth dependence of the projective charts on the parameter, locally in the manifold and parameter.

This is a conditional formulation, not a uniquely determined reading of the record. The source does not specify the projective category, the equivalence relation, the regularity of deformations, or whether actual deformations or only infinitesimal deformations are intended. It also does not say whether the question asks for the existence of some hyperbolic three-manifold with this property, asks about every manifold in a specified class, or seeks a criterion for a given manifold. Any required preservation of convexity, completeness, end geometry, cusp data, or boundary data remains unspecified.
\subsection{Short English statement}
The question asks whether the projective geometry induced by a hyperbolic metric on a three-manifold can vary nontrivially. A conventional real projective interpretation is given above, but the class of manifolds, the quantifier over them, and the permitted deformations require clarification. The supplied label is \texttt{partially\_solved}; this draft does not independently verify that assessment.
\subsection{Sources}
\begin{itemize}
\item Supplied local UnsolvedMath registry record, permanent ID \(30001310\), source problem code \texttt{catalogue entry}, title \emph{Deformations of Canonical Projective Structures on Hyperbolic Three-Manifolds}; collection: \emph{Oberwolfach Reports Open Problems}. No individual proposer is named; the supplied proposed year is \(2009\).
\item Supplied primary report citation: \emph{Differentialgeometrie im Gro\ss en (2010)}, Oberwolfach Reports, pp.\ 1943--1972, DOI \texttt{10.4171/owr/2009/35}. Source URL and supplied literature-source link labeled ``Differential geometry report'': \url{https://doi.org/10.4171/owr/2009/35}. The record describes the collection as drawn from mathematical workshops and their open-problem discussions.
\item Supplied literature-source link labeled ``Heusener--Porti rigidity theorem'': \url{https://doi.org/10.2140/gt.2011.15.2017}. No theorem statement or detailed bibliographic citation is included in the record.
\item Literature assessment dated \texttt{2026-08-21}: the supplied curation labels the problem partially solved and reports that canonical projective structures may be rigid or deformable depending on the manifold, that Heusener--Porti establish broad infinitesimal-rigidity and Dehn-filling results, and that no general classification is known. This is supplied dated curation derived from the OWR record, not an independently verified literature assessment; its status and cited claims require independent checking.
\item The registry labels the statement \texttt{exact} and reports that its mathematical content was checked against the official Oberwolfach report, with notation, context, and question form normalized without changing the source-level claim. This verification assertion is supplied metadata and has not been independently checked for this draft.
\end{itemize}
\end{document}
